% UNINTEGRATED EARLIER WHOLE-D18 APPENDIX, 25 October 2026.
% Existing general compact family; source preparation only.
% DRAFT: remaining internal proof bodies and review status are in the ledger.
% No worker import, enumeration, test, Lean or TeX execution.
\section{The earlier whole-\(D_{18}\) geometry}
\label{app:D18foundation}
This appendix supplies the head, prefix and fork inputs used before
the final dotted-head argument in the main text. We keep its geometric
steps in dependency order. In particular an old loss \(.61\) is used
only to establish head isolation; the later loss \(.58\) is proved
afresh. An acceptance record for an internal lemma does not replace
its proof.

Fix the original interior point and farthest vertex
\[
 z=E(p-x),\quad p^2=-1,\quad x\perp p,\quad \norm x<1.
\]
The entire incident diagram is \(D_{18}\), with actual unit normals
\[
 n_i=(e_i-e_{i+1})/\sqrt2\ (1\leq i\leq17),\qquad
 n_{18}=(e_{17}+e_{18})/\sqrt2 .
\]
For \(P_i=\sum_{j\leq i}e_j\), the duals are
\[
 \Omega_i=\sqrt2P_i\ (i\leq16),\quad
 \Omega_{17}=(P_{17}-e_{18})/\sqrt2,\quad
 \Omega_{18}=(P_{17}+e_{18})/\sqrt2.
\]
Thus \(x_1>\cdots>x_{17}>|x_{18}|\); the last coordinate is signed.
The fork duals have squares nine and mutual product eight, and
\(\ip{\Omega_i}{\Omega_f}=i\) for \(i\leq16\), \(f=17,18\).
All original reference arguments concern this one \(x\).

We use the main actual-point and payment lemmas.
A positive correction on actual normals retains every unknown original
wall sign by acuteness. Every core below has positive time; its negative
or null norm selects the future sheet. A spacelike corrected root is
used only through its positive finite projected point.
Every coefficient payment is made at its individual reference anchor
before its pressure is used. A zero-sum residual need not be orthogonal
to the reference; the general completed square retains that full cross.

\subsection{An initial window and actual head isolation}
\label{app:D18oldhead}
Put \(f_i^{\,o}=\sqrt i-.61/\sqrt2\).
The generic prefix proof in Appendix~\ref{app:prefixinventory}
applies through \(13\), with its actual \(M_{12}\) rank condition.
It gives a complete own/white list there and nonadjacent weak sites;
a weak cap is
\[
 c=Ap-\gamma\Omega_i-\Omega_1/2,\quad
 A^2=2i\gamma^2+2\gamma-\tfrac12,\quad
 \phi/2\leq\gamma<6/5,
\]
with own five, finite at least seven or dotted.
In particular a strong \(K=10,11\) exists before either new tail.
There are no spectators.

Deleting old \(1\) retains the full \(D_{17}\).
Finite classification leaves an isolated opposite
\(H=A_Hp-\eta\Omega_1\), \(A_H^2=2\eta^2-1\),
or one additional simple bond at old \(2\),
with \(A_H^2=2\eta^2+2\eta\).
A second contact makes a cycle; any other branch or four gives a
nonfinite principal at this rank.
In the isolated row owns zero, three and four have nonpositive time;
six has its actual \(6,3\) affine chain, and parallel has \(H+n_1\).
In the nonisolated row own zero has zero time and own three an
all-three triangle. Every remaining own is at least \(1/\sqrt2\),
and its actual edge has cosine \(2\eta+1\), so \(x_1>.643\).

We first extend the old-loss window through \(14,15\).
Deleting these retains \(A_{13}+D_4\) and \(A_{14}+D_3\).
There is at most one contact on each connected retained piece.
A connected full completion has rank eighteen, so is classical;
an exceptional finite diagram adds no bridge.
The target ratio is
\[
 T_o(q)=100/61-(q-.61)/[q(2q-.61)]<1.56
 \qquad(q=\sqrt{28},\sqrt{30}).
\]
Isolation and left-only contacts use the same Schur formulas as the
generic prefix. At isolated own five the left \(H_4\) finite point
has loss \(<.147+.174\); every finite own at least seven has loss
\[
 \frac{1.232}{\sqrt{28}+\sqrt{28-1.232}}+1/\sqrt{12}<.408;
\]
dotted loss is \(<.096+.5\).
The actual zero/small-own nulls are the double-fork affine \(D\)
chain, the four/simple star, the \(6,3\) chain and the parallel pair.
For a left-only own zero, far-three has negative time, next-far-three
and far-four have zero time; near-three/next-near-three give the
affine two-fork \(D\) principal and near-four gives affine \(B\).
Positive own-three with a left-three closes the intervening cycle.
All surviving left-only ratios other than the far helper exceed
\(1.59\), \(2\), or \(\sqrt{14}/2\).
These leave exactly the displayed pure helper.

The complete new right-hand data are
\[
\begin{array}{c|c|c}
i,\text{right contact}&\sigma&\Delta\\ \hline
14,\text{long leaf three}&\gamma+1/2&1/2\\
14,\text{either fork three}&\gamma+1/4&1/2\\
15,\text{middle three}&\gamma+1/2&1/2\\
15,\text{end three}&\gamma+1/4&5/8\\
15,\text{end four}&\gamma+1/(2\sqrt2)&1/4 .
\end{array}
\]
At \(15\) an end-five \(H_4\) is also allowed, with ratio \(>2.99\)
already at own zero.
The long/middle-three roles have own-zero double-fork affine
principals and own-three triangles; all remaining owns give ratio
at least \(1+1/\sqrt2\).
The fork/end-three roles have own-zero affine \(E_8\), own-three
cycles, own-four stars, own-six chains and the parallel null.
Their only extra old-loss cases are own five at either index and
finite own at least seven at \(15\).
The end-four own-zero kernel has coefficients
\((\sqrt2,2,2,1,1)\) along its actual cap, contacted fork,
center, other fork and old \(15\).
Bridges are only left-end-three joined to a \(D_4\) leaf, or
to the middle/end of \(D_3\), with residuals
\(1/28,1/30,19/120\). Own zero closes an actual cycle;
every positive own gives ratio \(>1.96\).
This exhausts the finite list, including both forks.

We exclude the extra right cases without assuming head isolation.
If old \(5\) is strong, a strong \(K=10,11\) reference has
energies \(>.823735,>.8472225\), sums \(<2.732,<2.886\),
terminal capacity \(>1/7\), and fifth capacity \(>.146\).
Every possible actual head supplies a separately funded pure gain
\(>.0135\): isolated five by its \(H_4\) point, finite at least seven
by \(H+n_1/\eta+2(\eta^{-1}-\eta)n_2\), dotted by \(H+n_1/\eta\),
and nonisolated by \(x_1>.643\).
The exact scalar seals for these rank-free points are proved in
Appendix~\ref{app:Aendpointcores} and the head calculation below;
their fifth prices are at most \(.1\), with zero sum on the first five.
If old \(5\) is weak, the actual weak-five point and strong \(6\)
give instead bare first-six energies \(>.8870,>.9104\),
with no additional head gain.
For a pure tail vector \(V\), equal to one through \(K\), with
\(\ip xV>G,\ V^2<R\), a legal \(1/7\) payment gives
\[
 \norm x^2>L_K+2(G-f_K^{\,o})/7-(R-K)/49.
\]
The extra head gain is included only in the already-strong-five row.

At own five set \(\alpha=\sqrt2(\phi/2+1/4)\);
\(\alpha^2=\phi+5/8>56/25\).
The positive actual left-\(H_4\) correction is
\((\phi/2)(3n_i/2+n_{i-1}+n_{i-2}/2)\);
its residual is \((5-3\phi)/8\).
The projected point has normalized displacement flat through \(i-3\).
It gives \(G>3.62,R=14\) at \(14\), and \(G>3.74,R=15\)
at \(15\). The least strong-five totals exceed \(1.009\),
and the bare weak-five totals are larger.
For every finite own at least seven at \(15\), put \(u=\gamma^{-2}\).
Use both actual positive cores
\[
 c+\gamma^{-1}n_{15}+2(\gamma^{-1}-\gamma)n_{14},
 \qquad
 c+\gamma^{-1}n_{15}+2(\gamma^{-1}-\gamma)n_{16}.
\]
Their square is \(-d,\ d=7-3u-4/u\), and their correction
cross is \(2-u\).
Averaging their spatial displacements gives
\(Y^2=A^2-3+u+2/u<A^2\).
Here \(.9<\gamma<1\), \(d\leq7-4\sqrt3<.0718\),
\(\alpha^2>2.645\), and
\(A^2=15\alpha^2-5/8\).
Thus \(Y/\alpha\) is flat through \(13\), with
\(G>3.677,R=15\).
The two strong-five totals exceed \(1.00519,1.00508\);
the bare weak-five totals exceed \(1.054\).
All coefficients are positive and selected products nonpositive,
so these are actual all-wall points or cores.

For adjacent hypothetical weak helpers with owns \(a,b<6/5\),
the actual product is at least
\(\tfrac12-ab+(a/b+b/a)/4\geq1-ab>-11/25\).
The native gap forces orthogonality.
Clearing the retained left path gives
\(I-t\mathbf1\mathbf1^{\mathsf T}\), \(t=(i-1)/(2i)\),
positive with positive inverse.
The two-helper mixed-point calculation of
Appendix~\ref{app:prefixinventory} uses only this matrix, times \(>2\)
and owns \(>4/5\); its weighted loss is
\(<7/16+1/(4\sqrt{10})<.61\).
The retained right block at \(14,15\) is actual \(D_3\).
Hence helpers are nonadjacent through \(15\).

We next prove old \(5\) strong; no head isolation is yet used.
If it is weak, \(6\) is strong and its actual point gives
\(x_1>.5965\). A finite weak-five own gives the stronger \(x_1>.65\);
its reference through \(1,6\) and a strong \(K=14,15\) has energy
\(>1.00845\), so the own must be dotted, \(1<a<6/5\).
With \(a_0=.5965\), the reference through \(1,6\), then every actual
strong knot, has initial energy \(>.7600299945\),
capacities \(P_1>.312168,\ P_6>.08806\), and terminal
\(P_{14}>.1361\).
All omitted weak knots are nonadjacent. Their complete worst debit is
\[
 [(.013586)^2+(.009296)^2+(.006871)^2+(.005344)^2]/2
 <.000174 .
\]
The underlying squared increments satisfy
\(\sum_{j=6}^{13}(\sqrt{j+1}-\sqrt j)^2>.2117655\)
and the next one \(>.0172465\), giving raw energies
\[
 q_{14}>.97162,\qquad q_{15}>.988867.
\]
These bounds follow directly by rational upper square bounds for
\(\sqrt{42},\ldots,\sqrt{210}\), or the increment identity
\(d_j^2=2j+1-2\sqrt{j(j+1)}\).
At \(a_0=.65\) the initial energy is \(>.7968668845\),
which proves the finite-own exclusion above.

For the actual full dotted weak-five foot put
\[
 U_5=\sqrt2aP_5-a^{-1}n_5+\Omega_1/2,\qquad
 T_5=\sqrt{10a^2+2a-\tfrac12}-\sqrt{1-a^{-2}}.
\]
Its pressure-minus-THIS-reference deficit is \(>.1203\).
To check the whole interval, the time derivative is \(<3.245\),
the coefficient of \(a\) in \(\ip w{U_5}\) is \(>2.4519\),
and the radical derivative is \(>1\), so the deficit decreases.
At \(a=6/5\) use time \(>4.03732\), radical \(<.552771\)
and \(\ip w{P_5}<1.733828\), giving \(>.120355\).
Pay \(\mu=.07,\ r=.088\), and
\(q=\mu(\sqrt2a+1/\sqrt2)-r\in(0,.081)\) at the separate
\(P_1,P_6\) anchors.
For \(Y_5=\mu U_5-qP_1-rP_6\),
\[
 Y_5^2=.0392a^2+.0049a^{-2}-.04928\sqrt2a+.02892<.005141 .
\]
Convexity reduces this bound to the two endpoints.
The general completed square, retaining its full reference cross,
gains \(>.011701\). If \(15\) is strong this already gives
\(\norm x^2>1.000568\).

If \(15\) is weak, \(14\) is strong and its pure helper
has time \(A_{15}^2=30g^2+2g-1/2\).
For finite own at least nine, the actual positive finite pair point
has spatial foot
\(g(\Omega_{14}+\Omega_{16})/2+\Omega_1/2\).
For all \(g>\cos(\pi/9)>.939\) its deficit exceeds
\[
 4.9529-.939(4.68151)-.42179=.13517211 .
\]
Indeed at \(.939\), the exact time and pair normalization give
\(H>4.9529\); the deficit increases since
\(H'>\sqrt{30}/1.066>5.13\).
Its tail residual has square \(<1\).
Paying \(.0962\) costs \(<.136049\) at \(P_{14}\) and
\(<.069\) at \(P_1\), no \(P_6\); its tail is orthogonal to \(Y_5\).
Its gain \(>.016752673964\), together with the previous one,
gives \(\norm x^2>1.00007367\).
For dotted \(15\), weakness forces \(1<g<10/9\);
the positive reciprocal core has tail square \(<1.28\) and
deficit \(>.16613\). Its \(.086\) payment fits the same two
anchors, has no \(P_6\) price and gains \(>.0191\).
The entire finite and dotted tails are covered by these derivative
and endpoint bounds, not by a finite sample.

For owns five, seven and eight do NOT insert the weak-five foot.
Their own tail points suffice on raw \(q_{14}\):
\[
\begin{array}{c|c|c}
\text{own}&\text{actual correction}&\text{deficit lower bound}\\ \hline
5&g n_{15}+(2/3)n_1+(1/3)n_2&.3776769\\
7&g n_{15}+n_1/2&.28605949\\
8&g n_{15}+n_1/2&.232552911 .
\end{array}
\]
The first row is the actual \(H_4\) finite point; the other two
are negative cores.
Their spatial feet are
\(g(\Omega_{14}+\Omega_{16})/2+\Omega_3/6\)
or \(g(\Omega_{14}+\Omega_{16})/2+\Omega_2/4\).
The native brackets \(.809<g_5<.81\),
\(.9009<g_7<.901\), \(.9238<g_8<.9239\),
time/radical pairs \(4.5554/.111,\ 5.0646/.249,\ 5.1913/.322\)
and \(\ip w{\Omega_3/6}<.2747,\ip w{\Omega_2/4}<.3115\)
give the table by direct subtraction.
The residual squares are \(<1\), including the initial variances
\(1/15,1/10\). Payment \(.09\) costs \(<.118\) at \(P_{14}\)
and \(<.009\) at \(P_6\), and gains \(>.0337\).
No weak-five foot is present, so there is no duplicate initial debit.
Thus old \(5\) is strong in every case.

A strong \(K=14,15\) now has the ordinary old-loss reference
energy \(>.90715\), first-five mean \(<.361\) and
fifth capacity \(>.146\).
A nonisolated head would give \(x_1>.643\), and its legal five-block
coordinate graft gains \(>(5/4)(.643-.361)^2>.0994\),
contradicting the norm budget.
The head is therefore ACTUALLY isolated, with own five,
finite at least seven, or dotted. This conclusion did not use any
later \(.58\) statement or center/fork hypothesis.

\subsection{The fresh \(.58\) window, with the two changed tail rows}
\label{app:D18freshwindow}
From now on put \(f_i=\sqrt i-.58/\sqrt2\).
The head isolation just proved remains valid; old \(.61\)-strength
is not silently called new strength.
At \(5\leq i\leq13\), deleting \(i\) retains
\(A_{i-1}+D_{18-i}\).
The complete white list is isolation; a left endpoint/\(D/E\)
three or end-four; right long-three or permitted fork-three;
or the full endpoint \(D_{18}\) bridge.
Every left-only own-four has an ACTUAL affine \(B\) principal:
coefficients \(\sqrt2\) on its cap, \(2\) on the old chain
\(i,\ldots,16\), and \(1\) on both forks.
Any extra left contact preserves all outside signs.
It is a forbidden positive future null, not a role declared strong
by an inadequate edge ratio.

The new ratio threshold is
\[
 T(q)=50/29-(q-.58)/[q(2q-.58)].
\]
It is increasing, with \(T<83/50\) for \(q<7\),
\(T<801/500\) for \(q<15/4\), and \(T<159/100\) for \(q<16/5\).
For near-three/own-five the ratio
\([i(\phi+1)-1]/\sqrt{2i(i+1)}\) increases with \(i\),
exceeds \(801/500\) at six and \(83/50\) at eight.
At five every own at least seven is stronger.
\(D\)-position ratio is at least \((i\phi/2+1)/\sqrt{2i}\),
already \(>159/100\) at five and \(>2927/1750\) from six.
\(E\), end-four, long-three and bridge rows are larger.
Right-fork-three at retained \(D_5,D_6,D_7\) has residual
at most \(3/8\) and numerator \(>\gamma+1/4\);
surviving \(\gamma\geq\phi/2\) gives ratio \(>83/50\).
The actual affine/time small-own boundaries are unchanged.

The new near-three/own-five at old five is absolutely impossible:
its actual right-\(H_4\) point has displacement \(Y\), with
\(\ip xY>18/5,\ Y^2<281/20,\ \sum Y_j<43/5\).
The independently rank-valid OLD strong \(11,12\) references
precede this new window.
Cauchy for \(Y+(3/4)P_K\) gives strict gaps
\(.0401,.78625625\).
Isolated five/seven/eight-or-larger have losses \(<.58\),
including the exact singleton maximum
\(\sqrt{10}+\sqrt2-4<.578\).
For isolated dotted \(i\geq6\), reciprocal-core loss is
\(<a_iu+\sqrt{u(1-u)}\), with \(a_i<37/250\);
its circle bound is \(<.5795\).
Only isolated dotted five remains in addition to the pure helper,
whose weak own range is again \(\phi/2\leq\gamma<6/5\).
The cutoff follows from \(R^2\geq4\gamma-2\);
it is not a cutoff on arbitrary fork opposites.

At fourteen/fifteen retain the COMPLETE white table above.
Now \(T(\sqrt{30})<1.64\).
Isolation, ordinary left-only rows, bridges, long/middle-three,
end-four and end-five are still strong by their displayed exact ratios.
Own-four is removed by the actual affine principal, not that comparison.
The previously excluded own-five and finite-end-three cases remain
absolute contradictions, independent of the numerical threshold.
The two genuinely new weak candidates are:
\[
 \begin{array}{c|c}
 14,\text{fork-three}&\gamma=\cos(\pi/7)\\
 15,\text{end-three}&1<\gamma<1.05 .
 \end{array}
\]
The second cutoff is
\(1.64\sqrt{5/8}-1/4<1.05\).
No formerly strong dotted interval has been discarded without rechecking.

Bulk nonadjacency supplies a strong \(K=11,12\) independently
of these two candidates. In the strong-five branch the raw energies
are \(>.8626,>.8843\), fifth capacity \(>.151\), and
terminal capacity \(>1/7\).
Every isolated head gives a pure separate gain \(>.0155\):
five has \(P_4/4\) pressure \(>.39307\) and \(.56\) payment,
with gain \(>.0155344\);
other heads have actual sum-one foot \(Z\), with
\(.8h-.16Z^2>.2893\), and \(.4\) payment at fifth price \(.08\),
gaining \(>.02914\).
In the weak-helper-five branch \(x_1>.5455\) and strong six
give initial energy \(>.74385\), stronger than the following
bare weak-five bounds.
For isolated dotted weak five put
\[
 U'=P_5+(u/2)(e_6-e_5),\quad .4<u<1,\quad \ip x{U'}>1.8185.
\]
Choose strong \(J=6,7\). Its zero-sum residual has norm at most
\(77/150,194/175\) and deficits \(>.119,>.221\).
For these EARLY references pay \(.19\) at \(J=6\) and \(.2\)
at \(J=7\). The first debit \(19/120\) fits \(>.15834\);
the old \(.21\) debit would not fit at \(K=11\).
The two actual gains exceed \(.02668,.04405\).
First-\(J\) energies \(>.69312,>.7139\) and exact straight tails
\((\sqrt K-\sqrt J)^2/(K-J)\) give, in every weak-five branch,
\[
 q_{11}>.8697,\qquad q_{12}>.8912 .
\]
No additional head gain is charged in those rows.

For the native-seven fourteen candidate use both actual side cores,
on side \(13\) or \(15\), as above.
Let \(u=\gamma^{-2}\), \(\alpha=\sqrt2(\gamma+1/4)\).
The complete rational seal is
\[
 1.2318<u<1.2322,\quad
 2.64914162<\alpha^2<2.649602,\quad
 d=7-3u-4/u<.05736<.2395^2 .
\]
Averaging the displacements with the FULL correction cross gives
\(Y^2=A^2-3+u+2/u\).
As \(A/\alpha>3.7163\), \(\sqrt d/\alpha<.1472\), and
\(3.5-u-2/u>.6438>.24\alpha^2\), the normalized foot is flat
through twelve, with \(G>3.569,\ R<13.76\).
The \(1/7\) tail payment gives bare weak-five totals
\(>1.0025,>1.0021\); strong-five plus its separate head gain
gives \(>1.0109,>1.0107\).

For the dotted fifteen candidate use the positive actual core
\(c+\gamma n_{15}\), with norm and cap product \(1-\gamma^2<0\),
own product zero, positive time and all other original signs termwise.
Here \(\alpha=\sqrt2(\gamma+1/4)\), \(A^2=15\alpha^2-5/8\).
Its normalized displacement is flat through fourteen and satisfies
\[
 A/\alpha>\sqrt{14.8}>3.847,\quad
 \sqrt{\gamma^2-1}/\alpha<.182,\quad R<15 .
\]
Thus \(G>3.665\), throughout \(1<\gamma<1.05\).
The same legal payments give bare totals \(>1.0046,>1.0042\)
or strong-five-plus-head totals \(>1.0130,>1.0128\).
Both candidates are impossible.
The fresh window through fifteen and helper nonadjacency are proved.

Consequently some \(K=14,15\) is strong.
The all-hole references now give effective energies
\(>.9231,>.94006\), sums \(f_K\), and terminal capacity \(>1/8\)
in all initial cases. The weak-five isolated point may now use
the later legal \(.21,.2\) payments; they are not added to the
earlier versions in the same norm budget.
For either genuine fork put \(V=P_{17}\pm e_{18}\), \(V^2=18\).
If its actual cosine were at least twelve, its pressure would exceed
\(3.902\); a \(1/8\) terminal payment gives
\[
 .9231+(3.902-3.332)/4-4/64=1.0031,\qquad
 .94006+(3.902-3.463)/4-3/64=1.002935.
\]
Thus BOTH actual fork cosines are less than twelve, without assuming
strong five, strong center sixteen or a sign for \(x_{18}\).

\subsection{The exact center-sixteen inventory}
\label{app:D18centerinventory}
Deleting actual \(16\) retains \(A_{15}\) and BOTH orthogonal
fork normals. Write an actual center opposite as
\[
 c=Ap-g\Omega_{16}-\eta\Omega_j-\delta\Omega_{17}
      -\epsilon\Omega_{18},
\]
omitting absent terms. Finite classification at the actual endpoint gives
exactly the following white possibilities:
\[
\begin{array}{c|c}
 \text{left contact}&\text{fork contacts}\\ \hline
 \text{none}&\text{none; one arbitrary finite label; or }(3,3),(3,4),(3,5)\\
 \text{endpoint three at }1,15&
 \text{none; one three/four; or both threes}\\
 \text{next-three at }2,14&\text{none}\\
 \text{endpoint four at }1,15&\text{none}.
\end{array}
\]
Two left contacts make a cycle. A fork added to a next-three
position has excessive arms \(1,2,13\); an internal four is not a
long finite classical tree. The full two-path join has rank eighteen,
beyond every exceptional type. The three-node fork list explicitly
retains \(H_3\), with labels \(3,5\). Thus no field restriction has
silently removed a white possibility.

The exact retained Gram gives
\[
 q_L=2\eta^2j(16-j)/16,\quad
 \Delta=1-q_L-\delta^2-\epsilon^2>0,\quad
 \sigma=g+\eta j/16+(\delta+\epsilon)/2,
\]
and \(A^2=32\sigma^2-\Delta\),
\(C_{16}=\sqrt{32}\sigma/\sqrt\Delta\).
The fresh target satisfies \(T(\sqrt{32})<1.641\), by positive
multiplication using \(5.656<\sqrt{32}<5.657\).

An isolated center has owns zero/three/four/six/parallel excluded
by time, the four-leaf-three star, the four/simple star, the \(6,3\)
chain and \(c+n_{16}\).
The left-\(H_4\) point at five gives
\(\ip x{\Omega_{16}}>\sqrt{32-g^{-2}}-.2>5.27\).
Finite owns at least seven have loss \(<.12+1/\sqrt{12}<.409\);
dotted owns have loss
\(<.09u+\sqrt{u(1-u)}<.5475\).
Thus every isolated center is strong.
For left-only own zero, the far-three time is negative, far-next-three
and far-four have zero time; the near-three/next-three affine \(D\)
and near-four affine \(B\) kernels use both actual forks.
Positive own three closes an all-three cycle, four gives the
four/simple star at \(16\), six the \(6,3\) chain and parallel a null.
All surviving endpoint-four, next-three and near-three ratios exceed
\(1.641\). The only no-fork weak role is
\[
 c=Ap-\sqrt2gP_{16}-e_1/\sqrt2,\quad
 a=A^2=32g^2+2g-\tfrac12,\quad
 \text{own five, finite }\geq7\text{ or dotted},\quad g<6/5 .
\]

Every left/fork bridge is strong or affine.
Endpoint-three plus one fork-three has residual \(9/32\);
own zero closes the long cycle, own three a triangle, and remaining
\(g\geq1/\sqrt2\) gives ratio \(>1.8\).
Endpoint-three plus one fork-four has residual \(1/32\)
and even its smallest own-zero ratio exceeds two.
Both fork-threes with the left contact have the same residual
and smallest numerator \(17/32\).
With no left contact, the pair \(3,3\) has own-zero four-cycle,
own-three triangle and all remaining ratios \(>1.7\).
The pair \(3,5\) already has own-zero ratio
\(\phi^3/2=\phi+1/2>2.1\).
The pair \(3,4\) has positive-own ratio \(>2.2\);
at own zero its actual reciprocal fork endpoint is the full \(B_{18}\)
positive projection, with cosine \(8+9\sqrt2/2>12\), impossible.

For exactly one fork with finite contact \(b\), own-three with
fork-three is affine. With fork-four/five its actual reciprocal
fork endpoint is \(D_{18}\), of cosine \(8+9b>12\).
Higher white labels give a direct ratio \(>1.8\).
Own four, six and parallel have the same actual star/chain/null,
using the UNCONTACTED fork in the star.
For \(g\geq\phi/2\), every fork label at least four has ratio at least
\(\phi/\sqrt2+1/2>23/14>1.641\).
For fork-three/own-five the actual left-\(H_4\) finite point gives
a normalized foot flat through thirteen, with \(G>3.85,R<16\).
The independent strong \(K=11,12\) references have energies
\(>.8626,>.8843\) and prices \(>1/7\), so the two legal norm
bounds exceed \(1.0299,1.0298\).
Thus the only positive-own nonpure weak role is single-fork-three,
with finite own at least seven or dotted and \(g<47/40\).
The cutoff follows from \(1.641\sqrt3/2-1/4<47/40\).

At own zero, single-fork-three/four/six have the actual \(E_8\),
four/simple and \(6,3\) nulls. Single-fork-five uses the actual
\(H_4\) on \(c,f,16,o\), with positive correction
\(b(3n_f/2+n_{16}+n_o/2)\).
Its spatial displacement is exactly \((b/\sqrt2)P_{15}\).
The finite point gives
\[
 S_{15}>\sqrt{2+8\phi}-(\sqrt5-2)>3.628 .
\]
This makes fifteen strong; upgrading its actual floor in the reference
gives
\[
 \norm x^2>.95678+2(.1313)(.165)+.165^2/15>1 .
\]
For every remaining single-fork own-zero cap, the fork length below
twelve leaves only finite white labels seven through ten.
Consequently the COMPLETE weak-center list is:
\[
 \begin{array}{c|c}
 \text{pure far-three}&\text{finite five/at least seven, or dotted}\\
 \text{one fork, own zero}&\text{white label }7,8,9,10\\
 \text{one fork-three, positive own}&\text{finite }\geq7
                         \text{ or dotted }1<g<47/40 .
 \end{array}
\]
Every boundary was dispatched before this necessary list was formed.

\subsection{Raw references, legal omissions and the center ellipse}
\label{app:D18rawreferences}
We give the reference construction independently of any selected
weak-center hypothesis. Whenever fifteen is strong, interpolate through
actual strong knots, omitting every weak knot by its neighbor chord.
Use the exact increment lower numerators, at denominator \(10^7\),
for \(d_i^2=(\sqrt{i+1}-\sqrt i)^2\), \(i=5,\ldots,16\):
\[
455488,385186,333704,294372,263340,238230,
217494,200080,185248,172465,161332,151549 .
\]
They follow from \(d_i^2(a-d_i^2)=1\), \(a=4i+2\), and
\(d_i^2>(a^2-1)/[a(a^2-2)]\).
Together with the initial energy \(>.6668157\), they give
\(E_{14}>.9241299,\ E_{15}>.9413764\).
The sum of all possible nonadjacent omission costs is \(<1/1024\).
With a native-five head, its actual first-seven relations force
five, six and seven strong; the sharper complete debit is
\[
 [(.0112)^2+(.008)^2+(.0061)^2+(.0049)^2]/2 ,
\]
giving \(>.94125107\).
These are squared-slope formulas for every allowed hole set,
not a numerical sample of selected weak sets.

The actual native-five head \(H_4\) foot is \(P_4/4\), with
pressure \(>.39307\), square \(1/4\).
Its \(.56\) zero-sum correction has fifth debit \(.112<.151\)
and gain \(>.0155344\).
For any other isolated head, its rank-free actual sum-one foot
has \(.8h-.16R>.2893\), and \(.4\) correction gives
gain \(>.02914\), at fifth debit \(.08<.151\).
Their complete scalar proof is the head-foot proof in
Appendix~\ref{app:B18references}; it concerns the identical actual
first three/four long roots and is independent of the terminal type.
When integrated, that shared proof belongs immediately after the generic
prefix's four weak-five pressures, rather than being repeated in both
classical terminal appendices.

The full endpoint-fifteen contracts are
\[
\begin{array}{c|c|c}
\text{initial/head case}&L_{15}\text{ lower}&\text{initial price}\\ \hline
\text{strong five, native-five head}&.95678547&.112\text{ at }P_5\\
\text{strong five, other head}&.96953&.08\text{ at }P_5\\
\text{weak helper five}&.97164&\text{actual }(1,6)\text{ reference}\\
\text{isolated dotted five, strong six}&.96525&.4/6\text{ at }P_6\\
\text{isolated dotted five, strong seven}&.96016&.4/7\text{ at }P_7 .
\end{array}
\]
The weak-helper row uses \(x_1>.5455\), initial energy \(>.74385\),
and strong six; its improvement over \(E_6<.7126\) is \(>.03125\).
For the isolated dotted rows, the BARE curved energies are
\(>.9212053,>.9035167\), first means \(<.3399,<.31938\),
and capacities \(>.1435,>.1366\).
They use only the head foot, gaining \(>.0440466,>.0566531\).
The isolated-five point \(U'\) is not also paid in these rows.
The native-five head is incompatible with weak isolated dotted five
by its actual local mutual in \((0,1/4)\).

Under the additional assumption that sixteen is weak and fifteen strong,
put
\[
 d=4-\sqrt{15}\in(.1270,.1271),\quad
 \Delta=S_{15}-f_{15}\geq0,\quad U=f_{16}-S_{16}\geq0 .
\]
The common reference has effective energy \(>.95678\)
and terminal capacity \(\beta\geq\sqrt{15}-\sqrt{14}>d\).
Retaining that unspent surplus in ONE square gives
\[
 2\beta\Delta+x_{16}^2+x_{17}^2+x_{18}^2<.04322,\quad
 x_{16}=d-U-\Delta .
\]
Hence
\[
 x_{17}^2+x_{18}^2<.04322-d^2+2dU-U^2 .
\]
For either signed \(t=x_{17}\pm x_{18}\), ordering gives
\(0<t<2d<.2541\), while a direct quadratic comparison gives
\(t<.233+1.1U\):
the difference of squares has constant \(>.000107\),
linear coefficient \(>.0042\), quadratic coefficient \(3.21\).
Thus, for \(P>0,z_0\geq0,z_0\leq10P\),
\begin{equation}\label{eq:D18centerellipse}
 PS_{16}+z_0(S_{16}+t)
 \leq f_{16}P+(f_{16}+.233)z_0 .
\end{equation}
This is a consequence of one funded original reference, not a new
independently optimized terminal energy.

\subsection{Pure weak centers: one cleared Gram and its full scalar dispatch}
\label{app:D18purecenter}
Suppose the center is weak and its cap is the pure far-three helper.
Own three/four/six/parallel were excluded in the inventory.
If fifteen were also weak, its actual cap is another pure helper.
The actual product is \(>1-gh>-11/25\), hence zero by the native gap.
Clearing actual \(A_{14}\) gives the positive two-cap Gram
\(\left(\begin{smallmatrix}8/15&-7/15\\-7/15&8/15\end{smallmatrix}\right)\);
both original forks remain orthogonal positive blocks.
The actual mixed-point loss is \(<7/16+1/(4\sqrt{10})<.58\),
contradicting both weak floors. Thus fifteen is strong and
\eqref{eq:D18centerellipse} is available.

Put
\[
 a=32g^2+2g-\tfrac12,\quad k=16g+\tfrac12,\quad
 P=k/(8\sqrt2),\quad r=17/32,\quad
 \bar c=c+\sum_{i=1}^{15}(16-i)n_i/16=Ap-PP_{16}.
\]
Then \(\bar c^2=r\) and \(k^2=8a+17/4\).
A finite own has \(.8<g<1\).
A dotted weak own has \(1<g<7/6\): the exact center cosine
is \((32g+1)/\sqrt{17}\), and at \(7/6\) its ratio already
exceeds \(821/500>T(\sqrt{32})\).
Set \(H=A-f_{16}P,\ K=3A-k\).
All intervals below have \(H,K>0\).

The seven actual fork white profiles and their lengths can be derived
before center strength is known.
For either assignment \(\{f,o\}=\{17,18\}\),
\(d=Bp-v\), own \(\delta=-\ip d{n_f}\geq0\), they are
\[
\begin{array}{c|c|c|c}
 &v&B^2&C_f\\ \hline
I&\delta\Omega_f&9\delta^2-1&3\delta\\
N&\delta\Omega_f+\Omega_o/2&9\delta^2+8\delta+5/4&(18\delta+8)/\sqrt{19}\\
N_4&\delta\Omega_f+\Omega_o/\sqrt2&9\delta^2+8\sqrt2\delta+7/2&9\sqrt2\delta+8\\
N_2&\delta\Omega_f+\Omega_{16}/2&9\delta^2+16\delta+7&9\delta+8\\
F&\delta\Omega_f+\Omega_1/2&9\delta^2+\delta-1/2&(18\delta+1)/\sqrt{19}\\
F_4&\delta\Omega_f+\Omega_1/\sqrt2&9\delta^2+\sqrt2\delta&9\sqrt2\delta+1\\
F_2&\delta\Omega_f+\Omega_2/2&9\delta^2+2\delta&9\delta+1 .
\end{array}
\]
Indeed deleting \(f\) leaves the full actual \(A_{17}\) path;
its one-node finite list is isolation, either-end three/four,
or either next-to-end three.
For a white \(\beta\Omega_j\), set
\(L=\ip{\Omega_f}{\Omega_j},J=\Omega_j^2\).
The edge formula is
\[
 C_f=(9\delta+\beta L)/(3\sqrt{\Delta_f}),\quad
 \Delta_f=1-\beta^2(J-L^2/9),
\]
and \((L,J)=(8,9),(16,32),(1,2),(2,4)\) gives the table.
Both signs are actual, with \(C_f<12\).
All own labels are initially retained.

Near-four forces \(\delta<1/2\), hence \(\delta=0\);
its acuteness square is \(-a/2-17/8<0\).
Near-next-three has \(\delta<4/9\) and square difference
\((\delta+1)[a(\delta-1)-(17/4)(\delta+1)]<0\).
Both are impossible.
For the other profiles clear only positive actual retained roots:
\[
\begin{array}{c|c|c|c}
 &v_0&s=\bar d^2&\rho=-\ip{\bar c}{\bar d}\\ \hline
I&0&1&\lambda\\
N&1/2&3/4&\lambda\\
F&1/16&17/32&\lambda+15/32\\
F_4&\sqrt2/16&1/16&\lambda+15/(16\sqrt2)\\
F_2&1/8&1/8&\lambda+7/16 .
\end{array}
\]
Here \(\lambda=-\ip cd\geq0\) is the ACTUAL mutual magnitude,
and
\(\bar d=Bp-[(\delta+v_0)P_{16}
+\delta(e_{17}\pm e_{18})]/\sqrt2\).
For \(N\), add \(n_o/2\).
For the far rows add the positive \(A_{15}\) inverse column.
The other fork is retained when cleared.
No projected normal is called an actual facet.

If \(\rho^2=rs\), \(\bar c+(\rho/s)\bar d\) is a forbidden
positive future feasible null.
If \(\rho^2>rs\), it is a negative core with square
\(r-\rho^2/s\), time \(A+\rho B/s>0\), selected products
that negative square and zero, and zero retained old products.
Every unknown wall is nonpositive termwise.
With \(t=x_{17}\pm x_{18}\), its pressure is contradicted when
\begin{equation}\label{eq:D18puregap}
 D=H+\frac{\rho}{s}
 \{B-v_0f_{16}/\sqrt2-(f_{16}+.233)\delta/\sqrt2\}
       -\sqrt{\rho^2/s-r}>0 ,
\end{equation}
provided \(z_0\leq10P_0\), where
\(P_0=P+\rho v_0/(s\sqrt2),\
z_0=\rho\delta/(s\sqrt2)\).
This funding condition precedes every use of the ellipse.

If \(\rho^2<rs\), use the positive actual finite projection
\[
 q_0=p+\alpha\bar c+\chi\bar d=(1+T)p-V,\quad
 \alpha=(sA+\rho B)/(rs-\rho^2),\
 \chi=(\rho A+rB)/(rs-\rho^2),\quad T=A\alpha+B\chi .
\]
All coefficients are positive, its norm is \(-(1+T)\), and its
normalization satisfies every original wall.
Its exact spatial direction and the SAME original maximum give
the subcritical estimates below. Both zero mutual and singular
boundaries are kept.

Here are the complete finite-own scalar constants.
They follow by positive squares, \(A'^2-32=17/a\), and
\(H'>.58\).
\[
\begin{array}{c|c|c}
g\text{ interval}&H\text{ lower}&K\text{ upper}\\ \hline
(.8,.85]&.4252&.695\\
[.85,.9]&.4542&.752\\
[.9,.95]&.4832&.809\\
[.95,1)&.5122&.864 .
\end{array}
\]
Also \(A/K>20/3\) throughout.
For dotted owns the corresponding complete table is
\[
\begin{array}{c|c|c|c}
g\text{ interval}&H\text{ lower}&K\text{ upper}&A/K\text{ lower}\\ \hline
(1,1.05]&.55235&.920&6.5\\
[1.05,1.1]&.58135&.974&6.5\\
[1.1,9/8]&.61035&1&6.5\\
[9/8,1.15]&.62485&1.028&6.46\\
[1.15,7/6)&.63935&1.05&6.4 .
\end{array}
\]
At a right endpoint the asserted \(K\) upper is proved by
\((k+K_{\rm upper})^2-9a\geq0\).
The five dotted gaps are
\(.0484,.037476,0,.045184,769/3600\);
the four finite gaps are
\(.012025,.005104,.027081,.008496\).
\(K\) increases, so these endpoint checks cover the intervals.
For \(A/K\), square \(k/A\), using \(k^2/a=8+17/(4a)\).
These are finite rational certificates for continuous ranges,
not sampled bounds.

For isolation, acuteness forces
\(\delta^2\geq a/(a-17/4)>1\).
Positive subcritical mutuals are exactly \(1/2,1/\sqrt2\).
Put \(\zeta=17/(4a)\) and
\(j=\zeta+k\lambda/(a\delta)\).
The finite projection gives
\[
 S_{16}+jt\geq\sqrt{16+2\zeta}-\sqrt{2(\zeta-j^2)},\quad j^2<\zeta .
\]
For finite \(g\), monotonicity of
\(A\sqrt{9\delta^2-1}-k\delta\) in both variables gives
\(\delta<8/5\), and the four strips have
\[
\begin{array}{c|c|c|c}
 &\zeta\text{ upper}&j-\zeta\text{ lower}&S_{16}\text{ lower}\\ \hline
1&.197&.180&3.611265\\
2&.175&.170&3.612025\\
3&.157&.161&3.61531\\
4&.141&.153&3.62223 .
\end{array}
\]
To reproduce each certificate, use \(t<.255\),
\(\sqrt{16+2\zeta}>4+\zeta/5\);
the resulting expression increases with \(j\) and decreases with
\(\zeta\) on that strip.
The four upper radicals are \(.332,.335,.335,.331\),
whose squared radicands are
\(.109742,.11195,.111752,.109128\).
For dotted \(g\), the same monotonicity gives
\(\delta<8/5,\ k/a>.41,\ \zeta<.127,\ j>\zeta+.128\);
the same rationalized test gives \(S_{16}>3.608275\).
All contradict weak \(S_{16}<3.59\).

For isolated supercritical mutuals use
\[
 3\delta-.172/\delta<B<3\delta-1/(6\delta),\quad
 \lambda<K\delta-A/(6\delta),\quad
 (f_{16}+.233)/\sqrt2<2.7033.
\]
Then \eqref{eq:D18puregap} exceeds
\[
 H+(.2967/K)\lambda^2+
 \{.2967A/(6K)-.172\}\lambda/\delta-\sqrt{\lambda^2-r}.
\]
For finite \(g\), the actual \(C_f<12\) gives \(\delta<4\)
and ensures the ellipse price; \(\lambda/\delta>.4\)
by the increasing positive-square test at \((g,\delta)=(.8,2)\).
Completing one square gives lower gaps
\(>.1292,.0931,.0593,.0295\) on the four strips.
For dotted \(g\), first take \(\lambda\leq4\).
Monotonicity gives \(\delta<5\), so
\(z_0<20/\sqrt2<10P\).
The ratio lower bounds \(.45,.55,.6,.7\), respectively on
\[
 (\lambda,g)>(\sqrt r,1),(1,1),(1.2,1),(1.4,1.15),
\]
follow by contradiction from
\[
 \lambda^2\leq \frac{t^2a}{9a-(k+t)^2}<\Lambda^2 .
\]
The last comparison reduces to the increasing quadratic
\(9a-(k+t)^2-(t^2/\Lambda^2)a>0\).
Its lower-bound signs are \(>.7586/r,>.6637,.715,1.89\).
With these ratios and the dotted strips, one square yields positive
lower gaps \(>.0156,.0046,.0075,.0004,.0058\).
The omitted small-\(\lambda\) portions are decreasing up to
\(1,1.2,1.4\); at those endpoints their gaps exceed
\(.26,.28,.153,.059\).
For \(\lambda\geq4\) no refined ellipse is needed:
\(t<2d<.2541\), \(K<1.081\), and the coarse core gap exceeds
\[
 .55235-.172(1.081)+.2599\lambda^2-\lambda>.52 .
\]
Thus every positive isolated mutual is excluded, including the entire
unbounded dotted alphabet. Only \(\lambda=0\) remains.

For near-three, acuteness is
\[
 Q(\delta,a)=(a-17/4)\delta^2-(17/4)\delta-3a/4-17/16\geq0 .
\]
Finite \(g\) gives \(\delta>1\), dotted \(g\) gives
\(\delta>39/40\); the latter follows from
\(Q(39/40,817/18)=-2693/19200<0\).
The actual edge gives \(\delta<23/9<3\).
The only positive subcritical mutual is \(1/2\).
Its full actual finite projection has \(T>300,j<1/2\) in
\(V=K_0(P_{16}+j(e_{17}\pm e_{18}))\);
hence the original maximum gives \(S_{16}+jt>3.76\).
For finite \(g\), use \(\delta<2,\ B/A<2,\ j<160/333\),
and \(T\)'s positive numerator \(>300(19/128)\).
For dotted \(g\), the sharper
\(13/10<\delta<3/2,\ B/A<1,\ j<1/3,\ T>8664/19\)
follows by the exact tests
\(A B-k(\delta+1/2)=1/2\) at \((g,\delta)=(1,3/2)\)
and below \(1/2\) at \((7/6,13/10)\).
Both contradict \(S_{16}+jt<3.59+.255/2=3.7175\).

Every supercritical near-three core is funded since
\(z_0-10P_0=(4\lambda/(3\sqrt2))(\delta-5)-10P<0\).
With \(J=k/2-4A/3\),
\[
 B>3\delta+4/3-19/(216\delta),\quad
 \lambda<K\delta-J,\quad J/K=A/(6K)-1/2 .
\]
In the finite strips the gap exceeds
\(.3238+(989/2160)\lambda^2+.3271\lambda
-\sqrt{4\lambda^2/3-r}\);
in the dotted strips it exceeds
\(.4292+(989/2625)\lambda^2+.3095\lambda
-\sqrt{4\lambda^2/3-r}\).
For \(\lambda\geq1\), completion of one square gives respectively
positive gaps \(>.1052,>.00406\).
Below one the displayed lower functions decrease and at one are
\(>.2127,>.2194\).
Thus near-three also forces mutual zero.
Its unique positive \(Q\)-root lies in \((1,3/2)\) for finite
\(g\), or \((39/40,33/32)\) for dotted \(g\).
Set \(\zeta=17/(4a)<.2\) or \(<.13\).
An isolated/near pair then has positive actual product
\(\delta_I[1/2-\zeta(\delta+1/2)]\);
two near caps have product
\(\delta-\zeta(\delta+1/2)^2>0\).
Two isolated caps have actual mutual
\(17/(4a-17)\in(0,1/4)\).
All caps are distinct by their original fork contacts.
Hence at least one fork must have a far profile.

For completeness here is the entire far-profile scalar certificate.
The exact positive-root sandwiches are
\[
\begin{array}{c|c|c}
 &M&M^2-B^2\\ \hline
F_2&3\delta+1/3&1/9\\
F_4&3\delta+\sqrt2/6&1/18\\
F&3\delta+1/6&19/36 .
\end{array}
\]
Always \(B<M\) and
\(B>M-E_1/\delta\), with
\(E_1=1/54,1/108,19/216\), respectively.
Both \(F_2,F_4\) are automatically supercritical, even at mutual zero.
The actual length gives
\(\delta<11/9<10(1/8)\) for \(F_2\), and
\(\delta<11/(9\sqrt2)<10(\sqrt2/16)\) for \(F_4\);
thus every such ellipse price fits.
Substitution in \eqref{eq:D18puregap} gives
\[
\begin{array}{c|c|c}
 &g<1&g>1\\ \hline
F_2&
 .2972+(989/360)\rho^2+.4906\rho-\sqrt{8\rho^2-r}&
 .3967+(1978/875)\rho^2+.4642\rho-\sqrt{8\rho^2-r}\\
F_4&
 .2972+(989/180)\rho^2+.6941\rho-\sqrt{16\rho^2-r}&
 .3967+(3956/875)\rho^2+.6568\rho-\sqrt{16\rho^2-r}.
\end{array}
\]
For positive mutual use \(\rho\geq15/16\) in \(F_2\);
for zero use exactly \(\rho=7/16\), whose radical is one.
For \(F_4\), \(\rho\geq15/(16\sqrt2)>5/8\).
The identity
\(c\rho^2-\sqrt{m\rho^2-r}\geq cr/m-m/(4c)\)
gives gaps \(>.0376,>.1852\) in the finite case and
\(>.032475,>.07256\) in the dotted case.
It retains every finite, parallel and dotted mutual.

For far-three with positive mutual \(\rho\geq31/32\).
The upper root gives \(\rho<K\delta\), and the actual
\(\delta<3\) funds \(z_0\leq10P_0\).
Positive-square monotonicity forces \(\delta>1\);
then the improved upper
\(B<3\delta+1/6-1/(12\delta)\) gives
\[
 D>H+c\rho^2+L\rho+C_1\rho/\delta-\sqrt{\rho^2/r-r},
 \qquad c=.2967/(rK).
\]
In finite strips, \(L>.0577,C_1>.1446,\rho/\delta>.48\);
in dotted strips \(L>.0546,C_1>.1322,\rho/\delta>.6\).
For the last ratio the exact monotone test at
\((g,\delta)=(1,155/96)\) has square gap \(37199/12288\).
For the finite ratio the corresponding test at \((.8,2)\)
has gap \(23399/102400\).
One square leaves respectively
\[
 H+.157621875/K-K/1.1868+.1252>.0047
\]
and, on the five dotted strips, gaps
\(>.0805,.0545,.0574,.0440,.0368\) with \(.1321\)
in place of \(.1252\).
Thus every positive mutual fails.

At zero far-three mutual the Schur determinant is \(1/16\).
The full actual \(A_{15}+n_o+c+d\) positive projection has
\[
 \alpha=(17A+15B)/2,\quad \chi=(15A+17B)/2,\quad
 T=[17(a+B^2)+30AB]/2 .
\]
Monotonicity of \(AB-k\delta-g-1/2\) gives
\(\delta<3/2,\ B/A<1\).
The spatial tail ratio is
\[
 j=\frac{(\chi/\alpha)\delta}
          {k/8+(\chi/\alpha)(\delta+1/16)}<1/2 .
\]
Its \(T>300\) for finite \(g\), \(T>1733/4\) for dotted \(g\).
The actual point therefore demands \(S_{16}+jt>3.76\)
against the same upper \(3.7175\).
This removes zero mutual too.

Every actual fork profile has now been covered under a weak PURE
center: near-four/next-three are impossible; isolated/near profiles
must be orthogonal and cannot form a pair; some far profile is
necessary, but every far profile fails.
All scalar interval checks above use explicitly displayed root
sandwiches, monotonicities and completed squares on complete ranges.
They do not posit a maximum native label or assume arithmeticity.
The pure weak-center branch is excluded.

\subsection{Nonpure centers and their actual neighbors}
\label{app:D18nonpurecenter}
Only the two nonpure weak roles from the complete inventory remain.
The own-zero/single-fork labels seven through ten will be excluded
directly at their actual neighbors; their older head-specific
refinements and separate strong-fifteen bootstrap are unnecessary.

First take single-fork-three with finite own \(.9<g<1\).
We do need fifteen strong for this branch.
If fifteen were weak, its actual pure helper has parameter
\(\phi/2\leq h<6/5\). Put \(y=h+1/30,\ t=g+1/4\);
the actual times satisfy
\[
 A_h^2=30y^2-8/15,\qquad A_c^2=32t^2-3/4,
 \qquad M=-\ip cd=A_hA_c-30yt .
\]
This retains the full fork contact.
It increases in both \(y,t\): positive derivative squaring gives
\[
 y^2A_c^2-t^2A_h^2=y^2(2t^2-3/4)+(8/15)t^2>0,
\]
and
\(1024t^2A_h^2-900y^2A_c^2
=t^2(1920y^2-8192/15)+675y^2>0\).
At the lower corner \(h=.8,g=.9\) the time/spatial square
gap is \(17.3085>0\); at \(h=6/5,g=47/40\) the comparison
with \(M=11/10\) has gap \(2861/8000>0\).
Thus the entire actual range has \(0<M<11/10\).
Native values leave \(M=1/2\) or \(M\geq1/\sqrt2\).

Positive clearing of \(A_{14}\) and the contacted actual fork gives
\[
 \bar d=A_hp-\sqrt2yP_{15},\quad \bar d^2=8/15,\qquad
 \bar c=A_cp-\sqrt2tP_{16},\quad \bar c^2=3/4,\quad
 \ip{\bar c}{\bar d}=-M .
\]
The other fork remains a positive retained unit block.
Both weak floors give deficits
\(D_h=A_h-\sqrt2y f_{15}>.43\),
\(D_c=A_c-\sqrt2t f_{16}>.60\).
At \(M=1/2\), the full positive actual projection has coefficients
\(43<\alpha<65,\ 37<\chi<56\), parameter \(T<959\);
its pressure would require
\(\alpha D_h+\chi D_c\leq\sqrt{1+T}-1<31\),
whereas the left side exceeds \(40.69\).
For \(M\geq1/\sqrt2\), the positive core
\(\bar d+(4M/3)\bar c\) has negative square
\(8/15-4M^2/3\).
Its pressure gap exceeds
\(.43+.8M-\sqrt{4M^2/3-8/15}\).
This decreases on the range and at \(11/10\) is \(>.27\).
All original wall signs hold by actual positive coefficients.
Thus fifteen is strong.

Use the one ending-fifteen reference of
Appendix~\ref{app:D18rawreferences}, with \(L>.95678\) and
unpaid terminal coefficient \(>.1313\).
For finite \(g\) the actual positive core is
\(c+n_{16}/g+n_{15}/(2g)\), of norm
\(-[1-3/(4g^2)]\).
Its cap/old-fifteen products vanish, its old-sixteen product is
negative, both fork products are negative, and time is positive.
With \(\alpha=\sqrt2(g+1/4)\), its normalized displacement \(V\)
has
\[
 V_j=1\ (j\leq14),\quad V_{15}=V_{16}=1-a,\quad
 V_{17}=2a+b,\quad V_{18}=\pm b,\quad
 a=\frac1{4g(g+1/4)},\quad b=\frac1{4(g+1/4)} .
\]
The full residual square is
\(1-2a+6a^2+4ab+2b^2<1.174\).
Its reference cross is \(-aw_{15}\), and its
pressure-minus-reference deficit is \(>.24441\):
use \(A_c/\alpha>3.964\),
\(\sqrt{1-3/(4g^2)}/\alpha<.28285\),
\(a>.2,\ w_{15}>.1313,\ f_{15}<3.463\).
A legal \(.13\) terminal payment therefore gives
\[
 \norm x^2>.95678+.26(.24441)-.0169(1.174)>1.000486.
\]
No absent \(P_{14}\) capacity or omitted negative cross is used.
This removes every finite own in the nonpure positive-own branch.

For dotted single-fork-three, \(g>1\), the actual center endpoint
has full forest \(A_{15}+A_2+A_1\).
For own zero and white \(\delta=\cos(\pi/m)\), \(7\leq m\leq10\),
the actual endpoint has \(A_{15}+I_2(m)+A_1\).
In both cases the positive original incident dual is \(P_{16}/4\):
\[
 q=Cp-\sqrt{C^2-1}P_{16}/4,\quad
 t=S_{16}/4<449/500,\quad E_q/E=C-\sqrt{C^2-1}\,t .
\]
These are actual full finite projections, so every unknown original
wall is satisfied. The reverse positive unit dual
\(u=(Cq-p)/\sqrt{C^2-1}\) lies in the dihedral block and satisfies
\[
 \ip yu=\frac{\beta-t}{1-\beta t},\qquad
 \beta=\sqrt{C^2-1}/C,\quad z=E_q(q-y).
\]
The point \(q\) is never asserted farthest; every original chamber
point is compared there with factor \(E/E_q\).

For dotted single-fork-three,
\(C=\sqrt{32}(g+1/4)/\sqrt{3/4}\), so
\(C^2>200/3>(204/25)^2\).
The increasing function \(C-(449/500)\sqrt{C^2-1}\) gives
\[
 E_q/E>4431/5000>7/8,\quad E/E_q<8/7,\quad
 \ip yu>5875/6824>43/50 .
\]
For own-zero/single-fork label seven through ten,
\[
 C=\frac{\sqrt8\,\delta}{\sqrt{1-\delta^2}},\quad
 C^2>648/19>(29/5)^2 ,
\]
as \(\delta>.9\). Hence
\[
 E_q/E>117/175>2/3,\quad E/E_q<3/2,\quad
 \ip yu>2900/3849>3/4 .
\]
The radical upper at \(29/5\) is \(40/7\);
the bounds \(\beta>197/200\) or \(124/125\) give the stated
reverse fractions by positive rational multiplication.
These two neighbor cases will use the following local facts.

\subsection{The two required central \(A_{15}\) facts}
\label{app:D18localA15}
Use zero-sum \(\mathbb R^{16}\), with
\[
 \omega_i=\sqrt2(P_i-iP_{16}/16),\quad
 D_i=i(16-i)/8,\quad C_{ij}=i(16-j)/8\ (i\leq j),\quad D_8=8 .
\]
At any actual \(A_{15}+A_2+A_1\) vertex, the actual central
edge has \(C_8>9/4\), without a farthest assumption.
At an actual \(A_{15}+I_2(m)+A_1\) vertex, \(7\leq m\leq10\),
with comparison at most \(3/2\), the \(A_{15}\) energy exceeds
\(7/16\). We prove both statements from the actual central opposite.

Write it \(d=Hq-v_A-v_S\), own \(\gamma\geq0\),
whole foreign cost \(Q=\norm{v_S}^2<1\).
Deleting eight leaves two \(A_7\) paths.
The finite white list is one three at
\(1,2,3,5,6,7\) or reflection; one four at \(1,7\)
or reflection; or one three on each path.
The third-from-end three is \(E_8\); the middle fourth position
is affine \(E_7\). No white five extends a full \(A_7\), and
a rank-fifteen join is classical with both new bonds simple.
The exact edge data are
\[
 \sigma=\gamma+
 \sum_{j<8}\beta_jj/8+\sum_{j>8}\beta_j(16-j)/8,\quad
 \Delta=1-Q-\sum_{j<8}\beta_j^2j(8-j)/4
             -\sum_{j>8}\beta_j^2(16-j)(j-8)/4 ,
\]
with \(H^2=8\sigma^2-\Delta,\ C_8=\sqrt8\sigma/\sqrt\Delta\).

We first remove own zero directly, without a foreign-shadow rule.
Two white paths close an actual affine cycle.
Single three at global \(3,13\) gives affine \(E_8\)
with arms \(1,2,5\); at \(4,\ldots,12\) use \(E_7\)
with arms \(1,3,3\).
White four at \(7,9\) has the four/simple star.
At endpoint four \(\Delta>0\) requires \(Q<1/8\),
so \(Q=0\) and \(H^2=-1/16\);
at next-three \(Q<1/4\) forces \(Q=0,H^2=-1/8\).
At endpoint three the retained bound is \(Q<9/16\),
but \(H^2=Q-17/32\).
In \(A_2+A_1\), the possible foreign contributions under that
bound are \(0,1/3,1/4,1/2\); simultaneous contact costs
at least \(1/3+1/4=7/12>9/16\).
Thus \(Q\leq1/2\), impossible.
For the high dihedral block, ANY positive contact costs
\[
 [4(1-\delta^2)]^{-1}>25/19>1 ,
\]
so there is no such contact at all. The remaining \(A_1\)
again gives only \(Q=0,1/4,1/2\) under the retained bound.
With no \(A\) contact the time is negative and its derivative
along that edge is zero.
Every affine kernel has positive actual coefficients and cap time,
selected products zero, all unknown products nonpositive;
it is a forbidden future feasible null.

For the first local statement, \(\gamma\geq\phi/2\) directly
gives \(C_8\geq\sqrt2\phi>9/4\).
Own four has a four/simple star at eight unless a white contact
hits \(7\) or \(9\); then
\(\sigma\geq1/\sqrt2+7/16,\Delta\leq9/16\), giving \(C_8>4\).
Own three has the actual \(E_7\) star at eight unless it contacts
one of the two three-node arms.
The remaining allowed contacts have
\((\sigma,\Delta)\geq/\leq
(13/16,1/16),(7/8,1/4),(15/16,9/16)\);
the weakest cosine is \(5\sqrt2/2>9/4\).
This proves the full central-edge statement.

For the high-dihedral second statement, all foreign cost is in
the remaining actual \(A_1\).
Own-three plus any white-three closes a cycle.
Pure far-three own-four/six/parallel have the actual star,
\(\gamma d+m_8+m_7/2\), and \(d+m_8\) nulls.
The remaining ordinary ratios satisfy
\[
\begin{array}{c|c|c}
\text{role}&\sigma\text{ lower}&\Delta\text{ upper before }Q\\ \hline
\text{near three}&.7+7/16&9/16\\
D\text{ three}&.7+1/8&1/4\\
E\text{ three}&.7+3/16&1/16\\
\text{white four}&1/2+1/(8\sqrt2)&1/8\\
\text{two threes}&1/2+1/8&1/8\\
\text{far three},Q>0&.8+1/16&5/16 .
\end{array}
\]
All exceed \(3/2\). Their actual edge demand is therefore
\(>(3\sqrt2-3/2)/\sqrt{17}>53/80\), giving energy
\(>2809/6400>7/16\).
Extra foreign cost only increases this demand.

For isolated \(A\) roles, own three/four/six/parallel have the
\(E_7\), star, chain and pair nulls.
Put \(\tau=(1-Q)/\gamma^2,\ H^2=\gamma^2(8-\tau)\).
For \(0<\tau<1\) the positive actual core
\(d+w_S+\gamma\tau m_8\) has spatial part
\(U=\gamma(\omega_8-\tau m_8)\), norm
\(-\gamma^2\tau(1-\tau)\), and \(U^2<H^2\).
If \(e_A\leq7/16\), Cauchy forces pressure loss divided by
\(\gamma\) greater than
\((1-\sqrt7/4)\sqrt{8-\tau}>5/6\),
against comparison loss at most \(3/4\).
At \(\tau=1\) it is a forbidden null.
For \(\tau>1\), the actual positive finite point gives
\[
 h_8\geq\sqrt{\frac7{8-\tau}}\,
              \{\sqrt7-\kappa\sqrt{\tau-1}\}.
\]
Foreign \(\sqrt{\tau-1}\leq2-\phi<.382\);
pure finite owns at least eight have
\(\sqrt{\tau-1}\leq\tan(\pi/8)<.415\).
The bracket is positive, and
\(h_8>13/5-(3/2)(.415)=791/400\), whose normalized value
is \(>2/3\), giving energy \(>4/9>7/16\).
Pure five uses the actual \(H_4\) finite point, with
\(\delta_5=(7-3\sqrt5)/16<1/36,\ H^2>4,\gamma>.8\):
\[
 h_8>\sqrt{8-\gamma^{-2}}
 -\kappa\sqrt{\delta_5}\sqrt{1+\delta_5/H^2}/\gamma
 >699/320 .
\]
Pure seven uses \(d+(4\gamma/3)m_8+(2\gamma/3)m_7\);
the native cubic brackets \(.9<\gamma<10/11\) give
\[
 h_8>\sqrt{8-\gamma^{-2}}
 -\kappa\sqrt{4/3-\gamma^{-2}}>103/50 .
\]
Both again give normalized height \(>2/3\).

For every pure far-three helper, without an upper own cutoff,
\[
 H^2=8\gamma^2+\gamma-17/32>8\gamma^2,\qquad \gamma>.8 .
\]
Five uses the actual correction
\(\gamma m_8+(2/3)m_e+(1/3)m_j\), residual
\(\delta=2/3-\gamma^2\in(0,1/36)\);
its finite point has \(U^2=H^2+\delta<(10/9)H^2\)
and pressure \(>H-.2525>(6/7)H\), giving energy \(>162/245\).
Finite owns at least seven use
\(d+\gamma m_8+2(1-\gamma^2)m_e\), negative square
\(-(1-\gamma^2)(4\gamma^2-3)\), radical at most \(1/4\).
Since \(H>2.5\), energy exceeds \(289/400\).
Dotted owns use \(d+m_8/\gamma\); its radical divided by
\(\gamma\) is at most \(1/2\), so the pressure ratio is
\(>1-3/(4\sqrt8)>2/3\), giving energy \(>4/9\).
All selected products and positive time are as displayed;
unknown wall signs hold termwise. This completes both local facts.

Return to the original actual neighbor.
In the dotted single-fork-three case, \(C_8>9/4\) and factor
\(<8/7\) give central normalized height
\(>31/(7\sqrt{65})\), hence \(e_{A_{15}}>961/3185\).
The orthogonal reverse \(A_2\) energy exceeds \(1849/2500\),
and their sum is
\(1+65813/1592500>1\).
In the own-zero/high-dihedral case, the second local fact gives
\(e_{A_{15}}>7/16\), while the reverse dihedral costs \(>9/16\).
The remaining actual \(A_1\) is nonnegative in each SAME neighbor
vector and is never charged twice.
Both weak-center families are excluded.

\subsection{Strong sixteen and the legal terminal append}
\label{app:D18strong16}
The complete center inventory, the entire pure-center exclusion,
and the finite/dotted single-fork-three and own-zero/high-dihedral
arguments cover every possible weak center.
Therefore \(S_{16}>f_{16}\) unconditionally.
No statement about strength of either fork was used.

Reassemble each endpoint-fifteen reference through actual strong sixteen,
then OMIT weak fifteen before applying any pressure.
Appending contributes \((4-\sqrt{15})^2\);
omitting fifteen costs
\((2\sqrt{15}-\sqrt{14}-4)^2/2<.00001\).
Since \(4-\sqrt{15}>.127\), the net gain is \(>.016119\).
The initial zero-sum corrections are supported through seven, so
their individual prices and full crosses are unchanged.
The endpoint-fifteen omission debit may conservatively be retained
in addition to this new debit; no fictitious floor is ever used.
The worst paid ending-sixteen effective energy is
\[
 L_{16}>.95678547+.127^2-.00001=.97290447>.9729 .
\]
The other four rows exceed
\(.985649,.987759,.981369,.976279\).
The unpaid \(P_{16}\) coefficient is its true last slope:
\(d_{15}\) if fifteen is strong or \((d_{14}+d_{15})/2\)
if fifteen is weak. It exceeds \(.127\) in every case.
All three late slopes are at least \(d_{15}\).

For either actual fork \(V=P_{17}\pm e_{18}\), the residual
\(V-P_{16}\) has square two and zero reference/initial cross.
If its actual cosine were at least ten, pressure would exceed
\(3.837\). A legal \(.12<\beta_{16}\) payment gives
\[
 \norm x^2>.9729+.24(3.837-3.59)-.0144(2)
 =1.00338>1 .
\]
Thus BOTH genuine fork cosines are less than ten for every head
and every original initial-five case.
The subsequent main-text uniform dotted-head argument strengthens
these data; it is not needed to prove them.

\subsection{The genuine fork inventory and the all-head near-pair bound}
\label{app:D18forkinventory}
Return to the seven actual fork profiles in
Appendix~\ref{app:D18purecenter}, now with \(C_f<10\).
For isolated or far-helper own three, an actual affine \(E_8\)
kernel has coefficients \(2\) on its cap, \(4\) on \(f\),
\(6\) on center \(16\), \(3\) on \(o\), and
\(5,4,3,2,1\) along \(15,\ldots,11\).
Own four has kernel \((\sqrt2,2,2,1,1)\)
on its cap, \(f,16,15,o\); own six has its \(6,3\) chain.
An extra far helper lies outside these principals.
All have positive cap time and all unknown wall signs termwise.
Parallel gives \(d+n_f\).
For near-three, own zero is isolated own three at the other fork;
own three is the actual four-cycle, own six the same chain.
Near-four now forces own zero and is the isolated own-four null
at the other fork; near-next-three likewise forces own zero and
is the four-leaf-three center star.
Far-four has no positive time at zero, and its only possible
positive own under length ten is three, already excluded.
These dispose of ALL formerly allowed near-four/next-three/far-four
values, not merely their finite members.

Isolated own five is exactly the actual \(H_4\) point already
computed in Appendix~\ref{app:D18centerinventory}.
Its displacement gives \(S_{15}>3.628\), hence fifteen is strong.
With terminal capacity \(>.1313\), the one-reference gain
\(2(.1313)(.165)+.165^2/15\) contradicts \(L_{15}>.95678\).
That point's computation did not use weak center sixteen.
The complete remainder is therefore
\[
\begin{array}{c|c}
I&\text{finite own }\geq7\text{ or dotted},\quad g<10/3\\
N&\text{own four/five, finite }\geq7\text{ or dotted},\
 g<(10\sqrt{19}-8)/18\\
F&\text{own five, finite }\geq7\text{ or dotted},\
 g<(10\sqrt{19}-1)/18\\
F_2&\text{own five or finite }\geq7,\quad g<1 .
\end{array}
\]
These are necessary roles, not examples.

The original ending-sixteen reference retains its terminal surplus.
For \(\Delta=S_{16}-f_{16}>0\),
\[
 \norm{x-w}^2<.0271-.254\Delta .
\]
Since \(w_{17}=w_{18}=0\), each genuine
\(t=x_{17}\pm x_{18}\) has \(t^2<.0542-.508\Delta\).
The square comparison
\((.233-\Delta)^2-(.0542-.508\Delta)
=.000089+.042\Delta+\Delta^2>0\)
gives
\[
 S_{17}+x_{18}<3.823,\qquad S_{17}-x_{18}<3.823.
\]
At cosine \(241/25=9.64\) the true fork demand has square
\(1944/133>3.823^2\); hence both fork cosines are \(<9.64\).

For a near-three or far-three fork cap, its actual other endpoint is
whole \(A_{18}\). Its original normalized dual height is
\(t<3.823/\sqrt{18}<.903\).
If its cosine \(C\geq6.82\), then
\[
 E_q/E>6.82-.903(6.747)=.727459>8/11,
 \qquad \ip y{u_{\rm reverse}}>.8 .
\]
The reverse positive dual is the new \(A_{18}\) endpoint dual.
Lemma~\ref{lem:localA18} applies at that actual \(q\), with factor
\(<11/8\) and endpoint height \(>.78\), and gives a contradiction.
That lemma is independent of the original head and never assumes
\(q\) farthest. Therefore
\[
 C_N,C_F<6.82,\qquad g_N<1.21,\qquad g_F<1.6 .
\]
These bounds follow from actual all-wall finite endpoints and
the same original maximum, not a whole-\(A_{18}\) farthest theorem.

Two near-three caps are impossible for EVERY head.
Write their owns \(g,h<1.21<3/2\) and times
\(A_g^2=9g^2+8g+5/4\).
Their spatial product is
\(K=8gh+(9/2)(g+h)+2\).
For \(P=gh,S=g+h\),
\[
 A_g^2A_h^2-K^2
 =17P^2-9(g^2+h^2)-(17/2)P-8S-39/16
 \leq(47/4)P-8S-39/16<0 ,
\]
because \(P<9/4,\ g^2+h^2\geq2P,\ P<3S/4,\ S<3\).
Their actual product is positive, violating acuteness.
They are distinct: coincidence would force \(g=h=1/2\),
whose actual cycle was already excluded.
This proof retains both fork signs and unequal owns.

\subsection{Finite heads at least nine, from their exact bare references}
\label{app:D18finiteheadslarge}
For a non-native-five head, the three BARE ending-sixteen references
are
\[
\begin{array}{c|c|c|c|c}
 &B&b\text{ upper}&P_J\text{ capacity}&J\\ \hline
\text{strong five}&.9565188&.3652&>.151&5\\
\text{isolated five, strong six}&.9373243&.3399&>.1435&6\\
\text{isolated five, strong seven}&.9196357&.31938&>.1366&7 .
\end{array}
\]
They are constructed by the same \(.016119\) append/omission
from bare fifteen energies
\(.9403998,.9212053,.9035167\).
No head or isolated-five pressure has yet been paid in these rows.
They remain flat first \(J\), have sum \(f_{16}\) and unspent
terminal price \(>.127\).

For an isolated finite head let \(u=\eta^{-2}\in(1,4/3)\).
The positive actual core
\[
 H+\eta u n_1+2\eta(u-1)n_2
\]
has square \(-\eta^2(7u-4-3u^2)\), zero head/first-root products,
negative second-root product, positive time and all unknown signs.
Its pressure is \(\ip xZ\geq h\), with
\[
 Z=(1-u/2)(e_1+e_2)+(u-1)e_3,\quad
 h=[\sqrt{2-u}-\sqrt{7u-4-3u^2}]/\sqrt2,\quad
 R=Z^2=3-4u+3u^2/2,\quad \sum Z_j=1 .
\]

If the native head label is at least ten, then
\(1<u\leq2-(2/5)\sqrt5<1.11\).
On this WHOLE interval \(F(u)=.8h-.16R\) decreases:
\((7u-4-3u^2)'>0\) and
\(F'<-.4/\sqrt2+.16<0\).
At \(1.11\), positive square bounds give
\(\sqrt{.89}>.94339,\sqrt{.0737}<.27149\);
thus \(h>.4748,\ R=.40815\), and \(F>.314536\).
Pay \(.4(Z-P_J/J)\) from the individual \(P_J\) anchor.
The general original square gives \(B+F-.8b+.16/J\),
which in the three rows exceeds
\(1.0108948,1.00660696,1.00152484\).
No old head lower gain has been subtracted from a paid lower bound.

In the fourth, weak-helper-five case the existing unheaded reference
has energy \(>.987759\), \(w_1=.5455\), \(w_2=w_3<.299\).
The exact head cross is
\[
 \ip wY< M(u)=(1.091-.2465u)/\sqrt2,\quad Y=\sqrt2Z,\qquad
 Q=Y-(\sqrt2/16)P_{16},\quad
 Q^2=T(u)=47/8-8u+3u^2 .
\]
Pay \(\mu\sqrt2/16<.018\) at the genuinely unspent terminal,
retaining this full cross. The norm lower bound is
\(.987759+2\mu(f-M)-\mu^2T\), \(f=\sqrt{2-u}-\sqrt{7u-4-3u^2}\).
All three functions decrease on the required interval.
Bounding \(f\) at the RIGHT and \(M,T\) at the LEFT endpoint gives
\[
\begin{array}{c|c|c|c|c|c}
u\text{ interval}&f>&M<&T\leq&\mu&\text{norm lower}\\ \hline
(1,1.05]&.7684&.5972&.875&.2&1.021239\\
[1.05,1.08]&.7125&.5885&.7825&.15&1.00735275\\
[1.08,1.1]&.6840&.5833&.7342&.14&1.00156468\\
[1.1,1.105]&.6778&.5798&.705&.14&1.001381\\
[1.105,1.11]&.6719&.5789&.698075&.13&1.0001415325 .
\end{array}
\]
The square brackets for the two radicals at the successive right
endpoints are
\((.9746,.2062),(.9591,.2466),(.9486,.2646),
(.9460,.2682),(.94339,.27149)\).
These prove every \(f\) entry; \(M\) uses \(\sqrt2>1.4142\)
and \(T\) is exact substitution.
The table is a complete interval certificate, not a sampled search.
All finite labels at least ten are excluded.

For native nine use its actual cubic:
\(8\eta^3-6\eta-1=0\) and
\(.93969<\eta<.9397\).
Exact products give \(1.1324<u<1.1325\).
Now
\(\sqrt{2-u}>.93139,\sqrt{7u-4-3u^2}<.28255\),
so \(h>.4587,\ R<.394\).
The strong-five bare row paid by \(.4\) gives
\(.9565188+.8(.4587-.3652)-.16(.394-.2)=1.0002788\).
For isolated-five rows set
\[
 Q_H=Z-P_J/J,\quad Q_5=U'-(5/J)P_J,\quad
 \ip{Q_H}{Q_5}=1-5/J .
\]
Their squares are \(R-1/J\) and at most \(77/150,194/175\).
Pay \(\mu=.4,\nu=.09,.1\), respectively.
The COMBINED anchor debits are \(.85/6<.1435,\ .9/7<.1366\).
The full one-vector energy bound
\[
 B+.8(h-b)+2\nu(1.8185-5b)
  -.16(R-1/J)-\nu^2Q_5^2-.8\nu(1-5/J)
\]
exceeds \(1.00125296,1.00128598\).
The positive mutual cross is explicitly subtracted.

A weak pure helper five cannot coexist with the native-nine head.
Its own is \(g=\phi/2\) or \(.9<g<1.2\), and its actual head
mutual is
\(\lambda=A_H\sqrt{10g^2+2g-1/2}-\eta(2g+1)\).
On the latter entire interval it increases, since the radical
derivative exceeds \(\sqrt{10}\) and \(A_H\sqrt{10}>2\eta\).
Native-nine brackets give
\(\lambda(.9)>.049875>0,\ \lambda(1.2)<.340629<1/2\),
a forbidden native gap.
At \(g=\phi/2\), direct positive squares give
\(\lambda<.8755(2.769)-.9396(2.618)<-.0356\),
contrary to acuteness.
All four original initial cases fail, excluding native nine.
The remaining finite heads are five, seven and eight.

\subsection{The all-head face envelope and an isolated-own cutoff}
\label{app:D18faceenvelope}
The paid ending-sixteen reference also controls EVERY asymmetric old
\(A_{16}\)-face direction. For \(|v|\leq1\), pay \(.127\)
at its unspent terminal toward \(Z=P_{17}+ve_{18}\).
The residual has square \(1+v^2\) and zero reference cross.
If \(\ip xZ\geq3.761+.0635v^2\), its original square exceeds
\[
 .9729+.254(.171)-.016129=1.000205
\]
because \(.254(.0635)=.016129\).
Thus
\begin{equation}\label{eq:D18faceenvelope}
 \ip x{P_{17}+ve_{18}}<3.761+.0635v^2,\qquad
 \frac{\ip x{P_{17}+ve_{18}}}{\sqrt{17+v^2}}<.9124 .
\end{equation}
For the second inequality the quotient of the upper bound decreases
with \(v^2\); its derivative numerator is
\(-1.602+.0635v^2<0\), and \(3.761^2<17(.9124)^2\).
The two old fork inequalities force the displacement direction of
ANY actual point on old \(A_{16}\) to have this form.
The original all-point maximum therefore gives cosine \(<10.94\)
for every such point.

Every isolated actual fork own \(g\) is less than \(15/8\).
Indeed its actual original fork endpoint is whole \(A_{17}+A_1\),
with cosine \(C=3g\), the new \(A_1\) being the cap itself.
At \(v=\pm1\), \eqref{eq:D18faceenvelope} gives original normalized
height \(<.90145\).
If \(g\geq15/8\), the actual endpoint \(q\) has
\[
 E_q/E>45/8-.90145(5.5355)=.635023525>5/8,\qquad
 E/E_q<8/5 .
\]
The reverse unit is exactly its whole \(A_1\) normal and has
height
\[
 >\frac{.984-.90145}{1-.984(.90145)}>.73.
\]
We prove the required local contradiction, rather than transporting
an original-farthest \(A_{17}\) theorem.

At an actual \(A_{17}+A_1\) vertex with factor \(<8/5\),
inspect central node eight in zero-sum \(\mathbb R^{18}\).
Here \(D=80/9,\ 2.98<\sqrt D<3,\sqrt{D-1}>2.8\).
Its retained paths are \(A_7,A_9\), with the finite white list
of Appendix~\ref{app:Acentralwindow}; explicitly single-three
positions \(1,2,3,5,6,7,9,10,16,17\), terminal fours
\(1,7,9,17\), two endpoint-threes or isolation.
The foreign whole \(A_1\) cost is \(Q=0\) or native \(Q\geq1/4\),
with \(Q<1\).
At own zero the endpoint-three time is \(Q-19/36\) with
\(Q<9/16\), allowing only \(0,1/4,1/2\), all impossible.
Next-three and endpoint-four require \(Q<1/4,Q<1/8\),
hence zero and negative time. Other zero-own positions have the
actual \(E_8/E_7\), four/simple or cycle nulls.
All are positive future actual combinations, feasible against
every unselected original wall.

For positive owns, ordinary white rows except native-five
single-white-three have actual cosine \(>5\).
The near-three residual is at most \(9/16\), with numerator
\(>.9+7/16\); \(D/E\) residual at most \(1/4\), numerator
\(>.9+1/20\); foreign far-three residual at most \(5/16\)
and the same numerator. A near-four with foreign cost satisfies
this latter bound. A pure near-four uses its ACTUAL mate at seven
or nine, of cosines
\((77+70\sqrt2)/(6\sqrt7)>10\) or \(9/2+4\sqrt2>10\).
White-four own-three at a far end has the affine \(F_4\) path;
near own-three/four has cosine \(>6\) by its ratio.
Two whites have residual at most \(9/80\), giving cosine \(>5\).
Own three/four/six/parallel at the other positions have their actual
cycle/star/chain/null boundary, with helper contacts outside
the selected principal. Hence every ordinary row has component energy
\[
 e_A>(5-8/5)^2/24=289/600>12/25 .
\]

For isolated own \(\gamma\), clear the whole foreign root and put
\(\tau=(1-Q)/\gamma^2\).
The same actual reciprocal core and positive finite point as in
Appendix~\ref{app:D18localA15} give
\[
 h_8>\sqrt{D-\tau}-\kappa\sqrt{\tau(1-\tau)}
 \quad(0<\tau<1),
\]
or
\[
 h_8\geq\sqrt{\frac{D-1}{D-\tau}}
 \{\sqrt{D-1}-\kappa\sqrt{\tau-1}\}\quad(\tau>1).
\]
The first case contradicts \(e_A\leq12/25\):
Cauchy would require loss \(>.3(2.8)=.84\), but the core costs
at most \(\kappa/2<.8\).
The null boundary is forbidden.
In the finite-point case, foreign radical \(<.4\) gives
\(h_8>2.16\); pure finite own at least eight has radical \(<.415\),
giving \(h_8>2.136\). Both normalized squares exceed \(12/25\).
Pure five uses the actual \(H_4\) finite projection,
\(\sqrt{\delta_5}<.1352,\gamma>.8,H/\gamma>2.7\),
and gives \(h_8>2.4296\).
Pure seven uses the actual two-root negative core, native
\(.9<\gamma<10/11\), and gives
\(h_8>2.76-(8/5)(.36)=2.184\).
Every selected product and future sign is the same direct actual-root
calculation already given; no native projection is assumed.

Pure far-three owns at least seven have the UNEQUAL exact times
\[
 H_1^2=D\gamma^2+(10/9)\gamma-19/36,\quad
 H_{17}^2=D\gamma^2+(8/9)\gamma-19/36 .
\]
Both exceed \(D\gamma^2\) for \(\gamma>.9\).
The finite core has radical at most \(1/4\) and \(H>8/3\),
giving normalized pressure \(>.85\).
The dotted core has radical divided by \(\gamma\) at most \(1/2\),
giving pressure ratio \(>1-\kappa/(2\sqrt D)>.7\).
Again energy exceeds \(12/25\).

The remaining native-five single-white-three row is paid jointly
with the known actual \(A_1\) reverse pressure.
Choose the \(H_4\) arm on the side AWAY from the helper.
Its positive finite point has full displacement \(U\), including
the foreign component, and residual \(\delta_5<.01825\).
The smaller cross \(C_{8j}\geq8/9\) gives \(H>2.45\).
Therefore the normalized point pressure exceeds
\(1-(8/5)(.1352)/2.45>.9\).
The actual foreign coefficient is \(<3/4\), so its correlation
with the unit \(A_1\) is \(<1/3\).
For the ONE vector
\[
 W=(387/800)u+(591/800)U/\norm U
\]
the positive pressures \(>.73,>.9\) give
\(\ip yW>1.0180125,\ W^2<1.0180125\);
hence \(\norm y^2>1\).
All other rows give the orthogonal sum
\(e_A+e_{A_1}>.48+.73^2>1\).
This proves the complete local contradiction and the original cutoff
\(g<15/8\), including every head and both signed fork choices.

\subsection{Two standalone local lemmas for a positive two-dual reverse}
\label{app:D18twodual}
The distant fork-pair exits below require a reverse direction which
can involve TWO positive duals. It cannot be replaced by a large
single endpoint coordinate.
We prove the two precise local lemmas here.

First let the whole actual forest be \(A_{17}+A_1\), comparison
\(\kappa<1.262\), and let
\[
 u=a e+b\omega_{17},\quad a,b>0,\quad
 a^2+(17/9)b^2=1,\qquad \ip yu>.8564 ,
\]
where \(e\) is the whole \(A_1\) unit.
Call only sites eight and nine weak if
\(h_i\leq\sqrt{D_i}-.8\).
If either is strong, its normalized dual has height \(>.73\)
and correlation with \(u\) less than \(1/4\)
(the exact bounds are \(\sqrt{8/(17\cdot10)},1/\sqrt{17}\)).
The single vector \(.75u+.55v_i\) has pressure \(>1.039\)
and square \(<1.07125\), giving norm \(>1.00675\).
Both sites must be weak.

Their finite-white list is the same full \(N=18\) list,
with all four retained side orders at least eight.
Thus zero-own endpoint-three has time \(Q-19/36\),
with native \(Q<9/16\) allowing only \(0,1/4,1/2\);
next-three has strict \(Q<2/m\leq1/4\);
endpoint-four has \(Q<1/8\).
They all have negative time, and every other zero-own role is
the actual positive affine kernel already displayed.
Using site seven instead would lose that strict side-order bound.

For positive ordinary ratios \(R\geq3/2\), the correct
comparison at the weaker \(\kappa=1.27\) and loss \(.8\)
reduces to the polynomial
\[
 -.21q^3+1.1729q^2-1.6q+.64>0
 \qquad(2.7\leq q\leq3.1).
\]
It is concave and its endpoint values are \(.737011,.695459\).
Near-four side orders at least eight also have ratio \(>3/2\).
Foreign exceptional helpers have ratio
\(>227/(25\sqrt{35})\), so cannot be weak.
Every isolated actual cap is strong by the same positive points;
their complete conservative losses at factor \(1.27\) are
\[
\begin{array}{c|c}
\text{dotted}&<(.193+1.285)/2=.739\\
\text{finite }\geq8&<5/26+1.27(.415)<.720\\
5\text{ finite point}&<25/76+1.27(1.01)/(6(.8))<.597\\
7\text{ negative core}&<10000/41877+1.27(.36)<.697\\
Q>0,\tau>1&<.193+1.27(.4)=.701 .
\end{array}
\]
These follow from the displayed \(\tau\) formulas, with
\(.193^2+1.27^2<1.285^2\);
the null boundary is unchanged.
Thus both weak sites must be pure far helpers with
\(\phi/2\leq\gamma<17/16\).
Their complete common/outward/reversed actual-product calculation
and positive two-cap Schur matrices were proved in
Appendix~\ref{app:Acentralwindow}, before its comparison.
In particular the reversed numerator has minimum
\(2763353243/50000000000>0\), unless both owns are five,
when its separate geometric distance-one obstruction applies.
The actual two-cap point has weighted mean loss
\[
 <.1+1.262(7/16)=.652125
 \quad\text{or}\quad <.1+1.262(3/10)=.4786,
\]
both \(<.8\).
Its entire actual \(A_1\) block is retained and every original wall
satisfied. This contradicts the two weak sites and proves the first
standalone two-dual lemma.

For the second, let the whole actual forest be \(A_{18}\), comparison
\(\kappa<1.27\), and
\[
 u=a\Omega_{18}+b\Omega_{17},\quad a,b>0,\quad
 \norm u=1,\quad L=a+2b\leq\sqrt{19/17}<1.058,\quad
 \ip yu>.85 .
\]
The inequality for \(L\) is Cauchy in the genuine \(N=19\)
inverse end block.
The complete own/white/null geometry in
Appendix~\ref{app:localAinputs} is independent of its later
single-endpoint consumer.
At sites five through fourteen, the same finite ratio polynomial
and isolated point losses just displayed are \(<.8\).
Hence possible weak roles are near-four at five/six or reflection,
or pure far helpers with \(\phi/2\leq\gamma<17/16\).
Their full actual helper nonadjacency gives weighted losses
\(<.655625,.481\), both below \(.8\), before near-four is removed.

A near-four actual mate has cosine \(>11\), hence normalized
height \(>7/8\). Its reverse correlation is
\(L\sqrt{2j/[19(19-j)]}\).
Left mates \(j=4,5\) have correlation \(<.206\) and their
two-direction sum contradicts the norm budget.
For a right mate \(j=14,15\), nonadjacency gives a left strong
\(k=5,6\), height \(>.703\).
The three full correlations are \(<.666,.234,.407\).
The ONE vector \(.5u+.6u_j+.4u_k\) has pressure \(>1.2312\),
square \(<1.45856\), and gives norm \(>1.00384\).
Thus every weak site is a nonadjacent pure helper.

Take symmetric rational heights \(H_{19-i}=H_i\), linear on both ends,
\[
 H_5=1.914,\ H_6=2.065,\ H_7=2.173,\ H_8=2.243,\
 H_9=H_{10}=2.277 .
\]
Their positive square tests \((H_i+.8)^2<D_i\) are the
earlier \(.87\)-window tests, with heights shifted by \(.07\).
Flatten a weak first/last block through six, and omit every remaining
weak knot by its chord.
The complete nonadjacent omission debit is
\(2(.043^2+.038^2+.036^2+.034^2)/4=.0028725\).
If fourteen is weak, the relevant last block has length
\(m=6,K=13,H_K=2.065\), energy \(E_0>.72555\) and
anchor \(>.11808\).
If fourteen is strong, \(m=5,K=14,H_K=1.914\),
energy \(E_0>.74793\), anchor \(.1159\).
The initial energy formulas are ordinary squared slope sums;
the worst two-weak-end value before omission is
\(2.065^2/6+.01772=.7284241666\ldots\);
the one-weak-end value is \(.7508121833\ldots\).
Stronger opposite-end patterns only increase these bounds.

Retain the exact reverse residual
\[
 G=u-(L/m)\Omega_K,\qquad
 \ip wG=0,\qquad G^2=1-2KL^2/(19m).
\]
Payment \(\mu=1/2\) spends only \(L/(2m)<.08817,.1058\)
at the appropriate actual strong anchor, fitting its own capacity.
One completed square gives
\[
 \norm y^2>E_0+.85-(H_K/m)L-1/4+KL^2/(38m).
\]
The gain decreases up to \(L=1.058\), where it is
\(>.299695\) or \(>.277477\).
Thus the norm exceeds \(1.025245\) or \(1.025407\).
The negative term in \(G^2\) and every individual coefficient
are retained. This proves the second standalone lemma.






